\documentclass[answers,12pt,addpoints]{exam}
\usepackage{import}

\import{C:/Users/prana/OneDrive/Desktop/MathNotes}{style.tex}

% Header
\newcommand{\name}{Pranav Tikkawar}
\newcommand{\course}{Stats 652}
\newcommand{\assignment}{Homework 1}
\author{\name}
\title{\course\ - \assignment}

\begin{document}
\maketitle

\begin{questions}

\question[0] \textbf{Problem 1.1.}

Define
\[
 g(\mu,\theta)=e^{-\mu^2\theta}+e^{-\theta}+\frac13\log\theta
\]
on
\[
 C=\{(\mu,\theta):0<\theta<\infty,\ \mu>\theta-1\}.
\]
Show that $g$ has only one critical point in $C$, and that it is a local maximum but
not a global maximum.

\begin{solution}
We first find all critical points of $g$ that lie in $C$.
The first-order partial derivatives are
\[
\frac{\partial g}{\partial\mu}
=
-2\mu\theta e^{-\mu^2\theta}
\]
and
\[
\frac{\partial g}{\partial\theta}
=
-\mu^2e^{-\mu^2\theta}
-e^{-\theta}
+\frac{1}{3\theta}.
\]

Because $\theta>0$ and $e^{-\mu^2\theta}>0$, the equation
\[
\frac{\partial g}{\partial\mu}=0
\]
holds if and only if
\[
\mu=0.
\]
Therefore every critical point of $g$ must have the form
$(0,\theta)$.

Substituting $\mu=0$ into the second score equation gives
\[
-e^{-\theta}+\frac{1}{3\theta}=0,
\]
or equivalently,
\[
e^\theta=3\theta.
\]

Define
\[
h(\theta)=e^\theta-3\theta,
\qquad \theta>0.
\]
Then
\[
h'(\theta)=e^\theta-3
\]
and
\[
h''(\theta)=e^\theta>0.
\]
Hence $h$ is strictly convex on $(0,\infty)$ and has a unique
minimum at
\[
\theta=\log 3.
\]
Moreover,
\[
h(0)=1>0,
\qquad
h(1)=e-3<0,
\]
so the intermediate value theorem implies that $h$ has a root
in $(0,1)$. Also,
\[
h(\log 3)=3-3\log 3<0,
\]
while $h(\theta)\to\infty$ as $\theta\to\infty$. Therefore $h$
has another root in $(\log 3,\infty)$. Since $h$ is strictly
convex and has only one minimum, these are its only two positive
roots.

However, a point of the form $(0,\theta)$ lies in $C$ only when
\[
0>\theta-1,
\]
which is equivalent to
\[
\theta<1.
\]
Thus the larger root does not give a point in $C$. Let
$\theta_0\in(0,1)$ denote the unique solution of
\[
e^{\theta_0}=3\theta_0.
\]
Then
\[
(0,\theta_0)
\]
is the unique critical point of $g$ in $C$.

We now classify this critical point. The relevant second
derivatives are
\[
g_{\mu\mu}(\mu,\theta)
=
-2\theta e^{-\mu^2\theta}
+
4\mu^2\theta^2e^{-\mu^2\theta},
\]
\[
g_{\mu\theta}(\mu,\theta)
=
-2\mu e^{-\mu^2\theta}
+
2\mu^3\theta e^{-\mu^2\theta},
\]
and
\[
g_{\theta\theta}(\mu,\theta)
=
\mu^4e^{-\mu^2\theta}
+
e^{-\theta}
-
\frac{1}{3\theta^2}.
\]

At $(0,\theta_0)$, these reduce to
\[
g_{\mu\mu}(0,\theta_0)=-2\theta_0<0,
\qquad
g_{\mu\theta}(0,\theta_0)=0.
\]
Because $\theta_0$ satisfies
\[
e^{-\theta_0}=\frac{1}{3\theta_0},
\]
we also have
\[
g_{\theta\theta}(0,\theta_0)
=
\frac{1}{3\theta_0}
-
\frac{1}{3\theta_0^2}
=
\frac{\theta_0-1}{3\theta_0^2}<0,
\]
where the final inequality follows from $\theta_0<1$.

Therefore the Hessian at the critical point is
\[
\nabla^2g(0,\theta_0)
=
\begin{pmatrix}
-2\theta_0 & 0 \\[4pt]
0 & \dfrac{\theta_0-1}{3\theta_0^2}
\end{pmatrix}.
\]
Both diagonal entries are negative, so the Hessian is negative
definite. Hence $(0,\theta_0)$ is a strict local maximum.

Finally, this local maximum is not a global maximum. Consider
the sequence
\[
(\mu,\theta)=(\theta,\theta),
\qquad \theta>0.
\]
Each such point lies in $C$, since
\[
\mu=\theta>\theta-1.
\]
Along this sequence,
\[
g(\theta,\theta)
=
e^{-\theta^3}
+
e^{-\theta}
+
\frac{1}{3}\log\theta.
\]
As $\theta\to\infty$, the first two terms tend to zero, whereas
\[
\frac{1}{3}\log\theta\to\infty.
\]
Consequently,
\[
g(\theta,\theta)\to\infty.
\]
Thus $g$ is unbounded above on $C$, and its unique critical point
is a strict local maximum but not a global maximum.
\end{solution}

\question[0] \textbf{Problem 1.2.}

Let $X_1,\ldots,X_n$ be a random sample from a
$\mathcal{N}(\mu,\sigma^2)$ distribution, where
\[
\mu\in\R
\qquad\text{and}\qquad
\sigma^2>0.
\]

\begin{parts}
\part Show that the likelihood goes to zero as
$(\mu,\sigma^2)$ approaches the boundary of the parameter space.

\part Show that the likelihood is not strictly log-concave in
$(\mu,\sigma^2)$.

\part Verify directly that the log likelihood is strictly concave
in $(\eta_1,\eta_2)$, where
\[
\eta_1=\frac{1}{\sigma^2},
\qquad
\eta_2=\frac{\mu}{\sigma^2}.
\]
\end{parts}

\begin{solution}
Let
\[
 v=\sigma^2,
 \qquad
 Q(\mu)=\sum_{i=1}^n(x_i-\mu)^2,
 \qquad
 S=\sum_{i=1}^n(x_i-\bar x)^2.
\]
Then
\[
 Q(\mu)=S+n(\mu-\bar x)^2,
\]
so the likelihood is
\[
 L(\mu,v)
 =(2\pi v)^{-n/2}
 \exp\left\{-\frac{Q(\mu)}{2v}\right\}.
\]

\begin{parts}
\part
Assume $S>0$, which holds with probability one under the normal model when $n\ge 2$. Since $Q(\mu)\ge S$,
\[
0\le L(\mu,v)
\le (2\pi v)^{-n/2}\exp\left\{-\frac{S}{2v}\right\}.
\]
The right-hand side tends to $0$ as $v\downarrow0$, uniformly in $\mu$. Also,
\[
L(\mu,v)\le (2\pi v)^{-n/2}\longrightarrow0
\qquad\text{as }v\to\infty.
\]
Finally, for fixed $\mu$, $L(\mu,v)$ is maximized over $v>0$ at $v=Q(\mu)/n$, and hence
\[
\sup_{v>0}L(\mu,v)
=
\left(\frac{n}{2\pi e\,Q(\mu)}\right)^{n/2}
\longrightarrow0
\qquad\text{as }|\mu|\to\infty,
\]
because $Q(\mu)=S+n(\mu-\bar x)^2\to\infty$. Therefore the likelihood tends to zero along every sequence leaving compact subsets of $\mathbb R\times(0,\infty)$. The assumption $S>0$ is necessary: if all observations equal $x_0$, then $L(x_0,v)=(2\pi v)^{-n/2}\to\infty$ as $v\downarrow0$.

\part
Ignoring constants that do not depend on $(\mu,v)$,
\[
\ell(\mu,v)
=
-\frac n2\log v-\frac{Q(\mu)}{2v}.
\]
Since
\[
Q'(\mu)=2n(\mu-\bar x),
\]
we have
\[
\frac{\partial\ell}{\partial\mu}
=
-\frac{n(\mu-\bar x)}{v},
\qquad
\frac{\partial\ell}{\partial v}
=
-\frac{n}{2v}+\frac{Q(\mu)}{2v^2}.
\]
The Hessian is
\[
H_{\mu,v}
=
\begin{pmatrix}
-\dfrac nv
&
\dfrac{n(\mu-\bar x)}{v^2}
\\[6pt]
\dfrac{n(\mu-\bar x)}{v^2}
&
\dfrac{n}{2v^2}-\dfrac{Q(\mu)}{v^3}
\end{pmatrix}.
\]
Its determinant is
\begin{align*}
\det H_{\mu,v}
&=
-\frac nv
\left(
\frac{n}{2v^2}-\frac{Q(\mu)}{v^3}
\right)
-\frac{n^2(\mu-\bar x)^2}{v^4}
\\
&=
\frac{n}{v^4}
\left[
Q(\mu)-n(\mu-\bar x)^2-\frac{nv}{2}
\right]
\\
&=
\frac{n}{v^4}
\left(S-\frac{nv}{2}\right).
\end{align*}
If
\[
v>\frac{2S}{n},
\]
then
\[
\det H_{\mu,v}<0.
\]
So the Hessian is indefinite at those points. Therefore the log
likelihood is not concave in $(\mu,v)$, and hence it is not strictly
log-concave in $(\mu,\sigma^2)$.

\part
Let
\[
\eta_1=\frac1v,
\qquad
\eta_2=\frac{\mu}{v}.
\]
Then
\[
v=\frac1{\eta_1},
\qquad
\mu=\frac{\eta_2}{\eta_1},
\]
and the parameter space is
\[
\{(\eta_1,\eta_2):\eta_1>0,\ \eta_2\in\R\}.
\]

Expanding the quadratic term gives
\[
\sum_{i=1}^n(x_i-\mu)^2
=
\sum_{i=1}^n x_i^2
-2\mu\sum_{i=1}^n x_i
+n\mu^2.
\]
Therefore,
\[
\ell(\eta_1,\eta_2)
=
C
+\frac n2\log\eta_1
-\frac{\eta_1}{2}\sum_{i=1}^n x_i^2
+\eta_2\sum_{i=1}^n x_i
-\frac{n\eta_2^2}{2\eta_1}.
\]

The first derivatives are
\[
\frac{\partial\ell}{\partial\eta_1}
=
\frac{n}{2\eta_1}
-\frac12\sum_{i=1}^n x_i^2
+\frac{n\eta_2^2}{2\eta_1^2},
\]
and
\[
\frac{\partial\ell}{\partial\eta_2}
=
\sum_{i=1}^n x_i-\frac{n\eta_2}{\eta_1}.
\]
Thus,
\[
H_\eta
=
\begin{pmatrix}
-\dfrac{n}{2\eta_1^2}
-\dfrac{n\eta_2^2}{\eta_1^3}
&
\dfrac{n\eta_2}{\eta_1^2}
\\[8pt]
\dfrac{n\eta_2}{\eta_1^2}
&
-\dfrac n{\eta_1}
\end{pmatrix}.
\]

The upper-left entry is negative:
\[
-\frac{n}{2\eta_1^2}
-\frac{n\eta_2^2}{\eta_1^3}<0.
\]
Also,
\begin{align*}
\det H_\eta
&=
\left(
-\frac{n}{2\eta_1^2}
-\frac{n\eta_2^2}{\eta_1^3}
\right)
\left(-\frac n{\eta_1}\right)
-
\left(\frac{n\eta_2}{\eta_1^2}\right)^2
\\
&=
\frac{n^2}{2\eta_1^3}>0.
\end{align*}
Therefore $H_\eta$ is negative definite for every
$\eta_1>0$. Hence the log likelihood is strictly concave in
$(\eta_1,\eta_2)$.
\end{parts}
\end{solution}

\question[0] \textbf{Problem 1.3.}

\begin{parts}
\part Show that if $f(x;\theta)$ is an exponential family in
$\theta=(\theta_1,\theta_2)$, then, for fixed $\theta_2$, it is an
exponential family in $\theta_1$.

\part Show whether, if every coordinate slice is an exponential family,
the joint family must be an exponential family. (Optional)
\end{parts}

\begin{solution}
\begin{parts}

\part
Suppose
\[
f(x;\theta_1,\theta_2)
=
\exp\left\{
\sum_{j=1}^k c_j(\theta_1,\theta_2)T_j(x)
+
d(\theta_1,\theta_2)
+
S(x)
\right\}I_A(x),
\]
where $A$ does not depend on $(\theta_1,\theta_2)$.

Fix $\theta_2=\theta_2^0$. Then
\[
f(x;\theta_1,\theta_2^0)
=
\exp\left\{
\sum_{j=1}^k c_j(\theta_1,\theta_2^0)T_j(x)
+
d(\theta_1,\theta_2^0)
+
S(x)
\right\}I_A(x).
\]
This is an exponential-family representation in $\theta_1$. Therefore,
for fixed $\theta_2$, the family is an exponential family in
$\theta_1$.

\part
No.

Let
\[
\mathcal X=A=\N,
\qquad
h(x)=2^{-x},
\qquad x\in\N.
\]
For each $j\in\N$, define
\[
U(x)=4^{-x},
\qquad
V(x)=3^{-x},
\qquad
T_j(x)=\mathbf 1\{x=j\}.
\]
Let the parameter space be $\Theta=\N^2$. For $(m,r)\in\Theta$, define
\[
Z_{m,r}
=
\sum_{x\in\N}h(x)
\exp\left\{mU(x)+rV(x)+\mathbf 1\{m=r\}T_m(x)\right\},
\]
and set
\[
f(x;m,r)
=
\exp\left\{
 mU(x)+rV(x)+\mathbf 1\{m=r\}T_m(x)
-\log Z_{m,r}+
\log h(x)
\right\}I_{\N}(x).
\]
Each $f(\cdot;m,r)$ is a pmf on $\N$, since $0<Z_{m,r}<\infty$.

We first record the following linear-independence fact, which we use below.
Suppose
\[
a_0+aU(x)+bV(x)+\sum_{j\in J}b_jT_j(x)=0
\qquad\text{for all }x\in\N,
\]
where $J$ is finite. If $x>\max J$, then all the $T_j(x)$ vanish, so
\[
a_0+a\,4^{-x}+b\,3^{-x}=0
\qquad\text{for all sufficiently large }x.
\]
Letting $x\to\infty$ gives $a_0=0$. Since $4^{-x}$ and $3^{-x}$ are linearly
independent as functions of $x$ (their ratio $(4/3)^{-x}$ is non-constant, so
no nontrivial combination can vanish at two or more values of $x$), this also
gives $a=b=0$. Evaluating at $x=j$ then shows $b_j=0$ for each $j\in J$.

Thus the functions
\[
1,\ U,\ V,\ T_1,T_2,\ldots
\]
are linearly independent in every finite combination.

Now fix $r=r_0$. Then
\[
f(x;m,r_0)
=
\exp\left\{
 mU(x)+\mathbf 1\{m=r_0\}T_{r_0}(x)
-\log Z_{m,r_0}+S_{r_0}(x)
\right\}I_{\N}(x),
\]
where
\[
S_{r_0}(x)=r_0V(x)+\log h(x).
\]
So for fixed $r_0$, this is an exponential family in $m$, with statistics
$U(x)$ and $T_{r_0}(x)$.

Likewise, if $m=m_0$ is fixed, then
\[
f(x;m_0,r)
=
\exp\left\{
 rV(x)+\mathbf 1\{r=m_0\}T_{m_0}(x)
-\log Z_{m_0,r}+S_{m_0}(x)
\right\}I_{\N}(x),
\]
where
\[
S_{m_0}(x)=m_0U(x)+\log h(x).
\]
So every slice with $m$ fixed is also an exponential family.

Suppose now that the full family were a $k$-dimensional exponential family for some
finite $k$. Then all log-density ratios would lie in one finite-dimensional vector
space spanned by $1$ and the fixed statistics of that representation.

Comparing the off-diagonal points $(2,4)$ and $(1,4)$ gives
\[
\log\frac{f(x;2,4)}{f(x;1,4)}=U(x)+\text{constant},
\]
so $U$ lies in that finite-dimensional space. Comparing $(4,2)$ and $(4,1)$ gives
\[
\log\frac{f(x;4,2)}{f(x;4,1)}=V(x)+\text{constant},
\]
so $V$ also lies in that space.

Finally, for each $j\in\N$,
\[
\log\frac{f(x;j,j)}{f(x;j,j+1)}=T_j(x)-V(x)+\text{constant}.
\]
Since $V$ is already in the space, it follows that every $T_j$ is also in the same
finite-dimensional space. But $T_1,T_2,\ldots$ are linearly independent, which is
impossible in a finite-dimensional space. This contradiction shows that the joint
family is not a finite-dimensional exponential family.
\end{parts}
\end{solution}

\question[0] \textbf{Problem 1.4.} (Optional)

Show that a $k$-dimensional exponential family is in minimal representation form if
and only if
\begin{parts}
\part $\sum_{j=1}^k a_jc_j(\theta)=a_0$ for all $\theta$ implies
$a_0=a_1=\cdots=a_k=0$;

\part $\sum_{j=1}^k b_jT_j(x)=b_0$ for all $x\in A$ implies
$b_0=b_1=\cdots=b_k=0$.
\end{parts}

\begin{solution}
Write
\[
f(x;\theta)
=
\exp\{c(\theta)^\top T(x)+d(\theta)+S(x)\}I_A(x).
\]

Suppose first that the representation is minimal. If
$a^\top c(\theta)=a_0$ for a nonzero vector $a$, choose an invertible linear
transformation whose last transformed coefficient is $a^\top c(\theta)$. That
coefficient is constant, so its term is independent of $\theta$ and can be absorbed
into $S(x)$. This removes one statistic and contradicts minimality. Hence (a) holds.

Likewise, if $b^\top T(x)=b_0$ for a nonzero vector $b$, choose an invertible linear
transformation whose last transformed statistic is $b^\top T(x)$. Since this
statistic is constant on $A$, its term can be absorbed into $d(\theta)$. This again
reduces the number of statistics, contradicting minimality. Hence (b) holds.

Conversely, assume (a) and (b). Fix $\theta_0\in\Theta$, $x_0\in A$, and define
\[
K(\theta,x)
=\log f(x;\theta)-\log f(x_0;\theta)-\log f(x;\theta_0)+\log f(x_0;\theta_0).
\]
This quantity is defined directly from the density $f$, so it does not depend
on which representation is used to compute it. Substituting
$f(x;\theta)=\exp\{c(\theta)^\top T(x)+d(\theta)+S(x)\}I_A(x)$, the terms
$d(\theta)$ and $S(x)$ cancel in the double difference, leaving
\[
K(\theta,x)=[c(\theta)-c(\theta_0)]^\top[T(x)-T(x_0)].
\]
Condition (a) implies that $\{c(\theta)-c(\theta_0):\theta\in\Theta\}$ spans
$\R^k$ (otherwise a nonzero vector orthogonal to this span would give a
forbidden affine relation among the $c_j$), and condition (b) similarly
implies $\{T(x)-T(x_0):x\in A\}$ spans $\R^k$. We may therefore choose
$\theta_1,\ldots,\theta_k\in\Theta$ and $x_1,\ldots,x_k\in A$ so that the
$k\times k$ matrix $[K(\theta_i,x_j)]_{i,j=1}^k$ has rank $k$.

Now suppose some representation of the same family used only $r<k$
statistics, say
$f(x;\theta)=\exp\{\sum_{l=1}^r c_l'(\theta)T_l'(x)+d'(\theta)+S'(x)\}I_A(x)$.
Because $K$ depends only on $f$, applying the same double-difference
identity to this representation gives
\[
K(\theta,x)=\sum_{l=1}^r\bigl[c_l'(\theta)-c_l'(\theta_0)\bigr]
\bigl[T_l'(x)-T_l'(x_0)\bigr].
\]
Evaluating at the chosen grid, $[K(\theta_i,x_j)]=\sum_{l=1}^ru_lv_l^\top$
where $u_l=(c_l'(\theta_i)-c_l'(\theta_0))_{i=1}^k$ and
$v_l=(T_l'(x_j)-T_l'(x_0))_{j=1}^k$. A sum of $r$ rank-one matrices has rank
at most $r$, since its column space is spanned by $u_1,\ldots,u_r$. But this
matrix was constructed to have rank $k$, so $r\ge k$. Hence no representation
with fewer than $k$ statistics exists, so the given representation is minimal.
\end{solution}

\question[0] \textbf{Problem 1.5.}

Let
\[
f(x;\mu)=
\frac{2}{\Gamma(1/4)}e^{-(x-\mu)^4},
\qquad
x,\mu\in\R.
\]
\begin{parts}
\part Show that this is a $k$-dimensional curved exponential family and specify $k$.

\part For a random sample of size $n$, determine the roots of the likelihood equation
and the MLE.

\part Find a minimal sufficient statistic for $\mu$.
\end{parts}

\begin{solution}
\begin{parts}
\part
Expanding the exponent,
\[
-(x-\mu)^4=-x^4+4\mu x^3-6\mu^2x^2+4\mu^3x-\mu^4.
\]
Thus
\[
f(x;\mu)
=
\exp\left\{
4\mu^3x-6\mu^2x^2+4\mu x^3
-\mu^4-x^4+\log\frac{2}{\Gamma(1/4)}
\right\}I_{\R}(x).
\]
Take
\[
T(x)=(x,x^2,x^3)^\top,
\qquad
c(\mu)=(4\mu^3,-6\mu^2,4\mu)^\top.
\]
The functions $1,x,x^2,x^3$ are linearly independent, as are the functions
$1,\mu,\mu^2,\mu^3$. Hence the representation is minimal and $k=3$. Since
$\{c(\mu):\mu\in\R\}$ is a one-dimensional nonlinear curve in $\R^3$ and does not
contain an open subset of $\R^3$, this is a curved exponential family.

\part
Ignoring constants,
\[
\ell(\mu)=-\sum_{i=1}^n(x_i-\mu)^4,
\qquad
\ell'(\mu)=4\sum_{i=1}^n(x_i-\mu)^3.
\]
Let
\[
F(\mu)=\sum_{i=1}^n(x_i-\mu)^3.
\]
Each summand is strictly decreasing in $\mu$, so $F$ is strictly decreasing. Also,
\[
F(\mu)\to\infty\quad(\mu\to-\infty),
\qquad
F(\mu)\to-\infty\quad(\mu\to\infty).
\]
Therefore the likelihood equation
\[
\sum_{i=1}^n(x_i-\mu)^3=0
\]
has exactly one root. Since $\ell(\mu)\to-\infty$ as $|\mu|\to\infty$, that root is
the unique global MLE.

\part
The factorization theorem gives sufficiency of
\[
T(X)=\left(\sum_{i=1}^nX_i,\ \sum_{i=1}^nX_i^2,\ \sum_{i=1}^nX_i^3\right).
\]
For sample points $x,y\in\R^n$, the log likelihood ratio is a polynomial in $\mu$.
It is independent of $\mu$ exactly when the coefficients of $\mu$, $\mu^2$, and
$\mu^3$ agree, which is equivalent to $T(x)=T(y)$. The likelihood-ratio criterion
therefore shows that $T$ is minimal sufficient.
\end{parts}
\end{solution}

\question[0] \textbf{Problem 1.6.}

Show that if $\eta$ represents the natural parameters of an exponential family, then
$\eta^*$ is also a natural parametrization if and only if
\[
\eta^*=A\eta+a
\]
for some nonsingular matrix $A$ and some $a\in\R^k$.

\begin{solution}
Throughout, ``natural parametrization'' is taken to mean a minimal
(full-rank) representation, as is standard; this guarantees $\eta$ and
$\eta^*$ have the same dimension $k$ (the family's minimal order) and lets
us invoke the characterization of minimality from Problem 1.4.

Assume both are minimal full $k$-dimensional natural representations of the same
family $\{f_\theta:\theta\in\Theta\}$, related to the original parameter $\theta$ by
bijections $\eta=\eta(\theta)$ and $\eta^*=\eta^*(\theta)$. Since both maps are
invertible, $\eta^*$ is in turn a well-defined function of $\eta$ alone, namely
$\eta^*(\eta):=\eta^*(\theta(\eta))$; this is the map we show to be affine. In the
first representation,
\[
f_\eta(x)
=
\exp\{\eta^\top T(x)+d_0(\eta)+S(x)\}I_A(x).
\]
By minimality, we can choose $x_0,x_1,\ldots,x_k\in A$ so that the matrix $D$ whose
$i$th row is $(T(x_i)-T(x_0))^\top$ is nonsingular. Writing out the log-density
ratio at each $x_i$ against $x_0$,
\[
\log\frac{f_\eta(x_i)}{f_\eta(x_0)}
=\eta^\top\bigl(T(x_i)-T(x_0)\bigr)+\bigl(S(x_i)-S(x_0)\bigr),
\]
the $d_0(\eta)$ terms cancel because they do not depend on $x$. Hence, setting
$b_i:=S(x_i)-S(x_0)$ (which does not depend on $\eta$) and
$b:=(b_1,\ldots,b_k)^\top$, the vector of log-density ratios has the form
\[
r(\eta)
=
\left(\log\frac{f_\eta(x_i)}{f_\eta(x_0)}\right)_{i=1}^k
=b+D\eta.
\]

Under the other minimal natural representation, with carrier function $S^*$, the
identical computation gives
\[
r(\eta)=b^*+D^*\eta^*(\eta),
\qquad
b^*_i:=S^*(x_i)-S^*(x_0),
\]
where $D^*$ is the (also nonsingular, by minimality of the second representation)
matrix with $i$th row $(T^*(x_i)-T^*(x_0))^\top$. Since $r(\eta)$ is the same
quantity computed from the same density $f_\theta$ in each representation, the two
expressions agree, so
\[
\eta^*(\eta)=(D^*)^{-1}D\eta+(D^*)^{-1}(b-b^*)=A\eta+a,
\]
and $A=(D^*)^{-1}D$ is nonsingular.

Conversely, suppose $\eta^*=A\eta+a$ with $A$ nonsingular. Then
\[
\eta=A^{-1}(\eta^*-a)
\]
and
\[
\eta^\top T(x)
=(\eta^*)^\top A^{-\top}T(x)-a^\top A^{-\top}T(x).
\]
Set
\[
T^*(x)=A^{-\top}T(x),
\qquad
S^*(x)=S(x)-a^\top A^{-\top}T(x).
\]
The second term is absorbed into the carrier $S^*$, while the first is
$(\eta^*)^\top T^*(x)$. Hence $\eta^*$ is also a natural parametrization.
\end{solution}

\question[0] \textbf{Problem 1.7.}

For a $k$-dimensional exponential family in natural form, prove the convexity,
derivative, Hessian, and MGF identities stated in the assignment.

\begin{solution}
Write the natural form as
\[
f_\eta(x)
=
\exp\{\eta^\top T(x)+d_0(\eta)+S(x)\}I_A(x).
\]
Let $h(x)=e^{S(x)}I_A(x)$ and define
\[
Z(\eta)=e^{-d_0(\eta)}
=\int_A e^{\eta^\top T(x)}h(x)\,d\nu(x).
\]
Assume the representation is minimal and work in the interior of the natural
parameter space, where differentiation under the integral is valid.

For any $a\ne0$,
\[
a^\top\nabla^2Z(\eta)a
=\int_A\bigl(a^\top T(x)\bigr)^2
 e^{\eta^\top T(x)}h(x)\,d\nu(x)>0.
\]
Minimality rules out $a^\top T(x)$ being almost surely constant. Thus
$Z(\eta)=e^{-d_0(\eta)}$ is strictly convex.

Now let
\[
A(\eta)=\log Z(\eta)=-d_0(\eta).
\]
Differentiating gives
\[
\nabla A(\eta)=\E_\eta[T(X)]
\]
and
\[
\nabla^2A(\eta)=\Cov_\eta(T(X)).
\]
For $a\ne0$,
\[
a^\top\nabla^2A(\eta)a
=\Var_\eta\bigl(a^\top T(X)\bigr)>0,
\]
so $A$ is strictly convex and $d_0=-A$ is strictly concave. In particular,
\[
\nabla d_0(\eta)=-\E_\eta[T(X)]
\]
and
\[
\frac{\partial^2d_0(\eta)}{\partial\eta_i\partial\eta_j}
=-\Cov_\eta(T_i(X),T_j(X)).
\]

For one observation, the data-dependent term $\eta^\top T(x)$ is linear in $\eta$.
Therefore the Hessian of the log likelihood is
\[
H_\ell(\eta)=\nabla^2\ell(\eta)=-\Cov_\eta(T(X)),
\]
so
\[
a^\top H_\ell(\eta)a
=-\Var_\eta\left(\sum_{j=1}^ka_jT_j(X)\right).
\]

Finally, whenever $\eta+t$ belongs to the natural parameter space,
\begin{align*}
M_{T,\eta}(t)
&=\E_\eta\left[e^{t^\top T(X)}\right]\\
&=e^{d_0(\eta)}
\int_Ae^{(\eta+t)^\top T(x)}h(x)\,d\nu(x)\\
&=\frac{Z(\eta+t)}{Z(\eta)}
=\exp\{d_0(\eta)-d_0(\eta+t)\}.
\end{align*}
\end{solution}

\question[0] \textbf{Problem 1.8.}

Find the MLEs, if they exist, in the following cases.
\begin{parts}
\part $X_1\sim\mathcal{N}(\mu,\sigma^2)$ and $X_2\sim\mathcal{N}(\mu,1)$ independently.
\part $X_i$ iid $\mathcal{N}(\theta,\theta^2)$.
\part $X_i$ iid $\mathcal{N}(\theta,\theta)$, $\theta>0$.
\part $X_i$ iid $\operatorname{Gamma}(\theta,\theta)$, $\theta>0$.
\part $X_i$ iid $\operatorname{Gamma}(\theta,1/\theta)$, $\theta>0$.
\end{parts}
The gamma density uses the shape--scale convention
\[
f(x;\alpha,\beta)=\{\Gamma(\alpha)\beta^\alpha\}^{-1}x^{\alpha-1}e^{-x/\beta}.
\]

\begin{solution}
\begin{parts}
\part
Set $\mu=x_1$ and let $\sigma^2\downarrow0$. The first normal density diverges while
the second remains positive and finite. The likelihood is unbounded, so no MLE
exists.

\part
The parameter space is $\R\setminus\{0\}$. Let
\[
A=\sum_{i=1}^n x_i^2,
\qquad
B=\sum_{i=1}^n x_i.
\]
Then
\[
\ell(\theta)=C-n\log|\theta|-\frac{A}{2\theta^2}+\frac{B}{\theta}-\frac n2.
\]
The score equation is
\[
n\theta^2+B\theta-A=0,
\]
with roots
\[
\theta_\pm=\frac{-B\pm\sqrt{B^2+4nA}}{2n}.
\]

Fix $A>0$; if $A=0$ then all $x_i=0$ and $\ell(\theta)\to\infty$ as
$\theta\to0$, so no MLE exists. On $(0,\infty)$ and on $(-\infty,0)$, as
$\theta\to0$ the term $-A/(2\theta^2)\to-\infty$ dominates
$-n\log|\theta|\to\infty$, so $\ell(\theta)\to-\infty$; as $|\theta|\to\infty$,
$-n\log|\theta|\to-\infty$ while the remaining variable terms vanish or stay
bounded, so again $\ell(\theta)\to-\infty$. Hence on each half-line $\ell$ is
continuous, tends to $-\infty$ at both ends, and so attains its supremum at an
interior critical point. Since $\theta_+\theta_-=-A/n<0$, exactly one root
lies in $(0,\infty)$ and one in $(-\infty,0)$, so $\theta_+$ and $\theta_-$
are the (unique, hence global) maximizers of $\ell$ on their respective
half-lines.

To compare the two, note that at any root of $n\theta^2+B\theta-A=0$ we have
$A/\theta^2=n+B/\theta$, so
\[
\ell(\theta)=C-n-n\log|\theta|+\frac{B}{2\theta}.
\]
Thus
\[
\ell(\theta_+)-\ell(\theta_-)
=-n\bigl(\log\theta_+-\log|\theta_-|\bigr)
+\frac B2\Bigl(\frac1{\theta_+}-\frac1{\theta_-}\Bigr).
\]
By Vieta's formulas, $\theta_++\theta_-=-B/n$, i.e.\ $\theta_+-|\theta_-|=-B/n$,
so $\theta_+<|\theta_-|$ exactly when $B>0$, making the first term positive
exactly when $B>0$ (negative when $B<0$, zero when $B=0$). Also
$\theta_+\theta_-=-A/n<0$ gives
$\tfrac1{\theta_+}-\tfrac1{\theta_-}=\tfrac{\theta_--\theta_+}{\theta_+\theta_-}>0$
always, so the second term shares the sign of $B$ as well. Hence
\[
\ell(\theta_+)>\ell(\theta_-)\iff B>0,
\qquad
\ell(\theta_+)=\ell(\theta_-)\iff B=0,
\]
so the global MLE is $\theta_+$ when $B>0$, $\theta_-$ when $B<0$, and both
roots tie for the global maximum when $B=0$.

\part
Here
\[
\ell(\theta)=C-\frac n2\log\theta-\frac{A}{2\theta}+B-\frac{n\theta}{2}.
\]
The score equation is
\[
n\theta^2+n\theta-A=0.
\]
If $A>0$, the unique positive root is
\[
\widehat\theta=\frac{-1+\sqrt{1+4A/n}}{2},
\]
and it is the unique MLE. If $A=0$, the likelihood is unbounded as
$\theta\downarrow0$, so no MLE exists.

\part
Let
\[
S=\sum_{i=1}^n x_i,
\qquad
G=\sum_{i=1}^n\log x_i.
\]
Then
\[
\ell(\theta)=C+(\theta-1)G-\frac{S}{\theta}-n\log\Gamma(\theta)-n\theta\log\theta.
\]
Writing $\psi$ and $\psi_1$ for the digamma and trigamma functions, the score and
its derivative are
\[
U(\theta)=G+\frac{S}{\theta^2}-n\psi(\theta)-n(\log\theta+1)
\]
and
\[
U'(\theta)=-\frac{2S}{\theta^3}-n\psi_1(\theta)-\frac n\theta<0.
\]
Moreover, $U(\theta)\to\infty$ as $\theta\downarrow0$ and
$U(\theta)\to-\infty$ as $\theta\to\infty$. Hence the MLE is the unique positive
solution of $U(\theta)=0$.

\part
Now
\[
\ell(\theta)=C+n\theta\log\theta-n\log\Gamma(\theta)+(\theta-1)G-\theta S.
\]
Thus
\[
U(\theta)=n(\log\theta+1)-n\psi(\theta)+G-S
\]
and
\[
U'(\theta)=n\left(\frac1\theta-\psi_1(\theta)\right)<0.
\]
Also,
\[
\lim_{\theta\to\infty}U(\theta)
=n+G-S
=\sum_{i=1}^n(1+\log x_i-x_i)\le0.
\]
The inequality is strict unless every $x_i=1$. If the data are not all equal to $1$,
there is a unique positive root of $U(\theta)=0$, and it is the MLE. If every
$x_i=1$, there is no finite MLE; the likelihood approaches its supremum as
$\theta\to\infty$.
\end{parts}
\end{solution}

\question[0] \textbf{Problem 1.9.}

Suppose $X_1,\ldots,X_n$ are iid Bernoulli$(\theta)$. For an integer $k$, find the
UMVUE of $\theta^k$ and determine when it does not exist.

\begin{solution}
Let
\[
T=\sum_{i=1}^nX_i\sim\operatorname{Binomial}(n,\theta).
\]
Then $T$ is complete and sufficient. For $0\le k\le n$, using the falling factorial
$(t)_k=t(t-1)\cdots(t-k+1)$,
\[
\E_\theta[(T)_k]=(n)_k\theta^k.
\]
Therefore
\[
\widehat{\theta^k}
=\frac{(T)_k}{(n)_k}
=\frac{\binom{T}{k}}{\binom{n}{k}}
\]
is unbiased and is the UMVUE by the Lehmann--Scheff\'e theorem.

Every expectation of a statistic on the finite Bernoulli sample space is a polynomial
in $\theta$ of degree at most $n$. If $k>n$, the function $\theta^k$ has degree greater
than $n$; if $k<0$, it is not a polynomial. Hence no unbiased estimator exists in
those cases. The UMVUE exists exactly for
\[
k\in\{0,1,\ldots,n\}.
\]
For $k=0$, the estimator is the constant $1$.
\end{solution}

\question[0] \textbf{Problem 1.10.}

Prove that any function of a complete statistic is complete.

\begin{solution}
Let $T$ be complete and set $V=g(T)$. Suppose $a$ is measurable,
$\E_\theta|a(V)|<\infty$, and
\[
\E_\theta[a(V)]=0
\]
for every $\theta$. Define $b(t)=a(g(t))$. Then
\[
\E_\theta[b(T)]=\E_\theta[a(g(T))]=0
\]
for every $\theta$. Since $T$ is complete, $b(T)=0$ almost surely for every $\theta$.
But $b(T)=a(V)$, so $a(V)=0$ almost surely for every $\theta$. Therefore $V$ is
complete.
\end{solution}

\question[0] \textbf{Problem 1.11.}

Let $X_1,\ldots,X_n$ be iid Poisson$(\lambda)$.
\begin{parts}
\part Show that no unbiased estimator of $|\lambda-1|$ exists.
\part For $k\in\R$, find the UMVUE of $\lambda^k$, if it exists.
\part For $a>0$, find the UMVUE of $e^{-a\lambda}$. Explain the problem when
$n<a$, especially when $a=2n$.
\end{parts}

\begin{solution}
Let
\[
T=\sum_{i=1}^nX_i\sim\operatorname{Poisson}(n\lambda).
\]
The statistic $T$ is complete and sufficient.

\begin{parts}
\part
Suppose an integrable unbiased estimator existed. Rao--Blackwellization would
give a function $h(T)$ satisfying $\E_\lambda[h(T)]=|\lambda-1|$. Then
\[
e^{n\lambda}\E_\lambda[h(T)]
=\sum_{t=0}^{\infty}h(t)\frac{(n\lambda)^t}{t!}.
\]
Because $h(T)$ is integrable for every $\lambda>0$, this power series converges
absolutely on every bounded interval and defines an analytic function of $\lambda$.
It would have to equal
\[
e^{n\lambda}|\lambda-1|,
\]
which is not differentiable at $\lambda=1$. This is impossible. Therefore no unbiased
estimator exists.

\part
If $k\in\{0,1,2,\ldots\}$, then
\[
\E_\lambda[(T)_k]=(n\lambda)^k,
\]
so
\[
\widehat{\lambda^k}=\frac{(T)_k}{n^k}
\]
is unbiased and is the UMVUE. For any other real $k$, the function
$e^{n\lambda}\lambda^k$ does not have an ordinary power-series expansion at
$\lambda=0$. Since the expectation of any integrable function of $T$ produces such a
power series, no integrable unbiased estimator exists.

\part
The pgf of $T$ gives
\[
\E_\lambda[r^T]=e^{n\lambda(r-1)}.
\]
Taking $r=1-a/n$ yields
\[
\widehat{e^{-a\lambda}}
=\left(1-\frac an\right)^T.
\]
It is unbiased and, as a function of the complete sufficient statistic, is the UMVUE.
When $a>n$, the base is negative, so the estimator changes sign even though the
target is positive. When $a=2n$, the estimator is $(-1)^T$ and
\[
\Var_\lambda((-1)^T)=1-e^{-4n\lambda}.
\]
Thus the estimator remains unbiased but is highly variable and oscillates between
$-1$ and $1$.
\end{parts}
\end{solution}

\question[0] \textbf{Problem 1.12.}

For a regular one-parameter family and a one-to-one reparametrization
$\gamma=g(\theta)$, show that
\[
I(\gamma;X)=\frac{I(\theta;X)}{\{g'(\theta)\}^2}.
\]

\begin{solution}
By the chain rule,
\[
U_\gamma(X)
=\frac{\partial}{\partial\gamma}\log f(X;\theta(\gamma))
=U_\theta(X)\frac{d\theta}{d\gamma}
=\frac{U_\theta(X)}{g'(\theta)}.
\]
Squaring and taking expectations gives
\[
I(\gamma;X)
=\E[U_\gamma(X)^2]
=\frac{I(\theta;X)}{\{g'(\theta)\}^2}.
\]
\end{solution}

\question[0] \textbf{Problem 1.13.}

For a one-dimensional exponential family in natural form, prove
\[
I(\eta;X)=-d_0''(\eta).
\]
In the original parameterization, prove
\[
I(\theta;X)
=-d''(\theta)+\frac{d'(\theta)c''(\theta)}{c'(\theta)}.
\]

\begin{solution}
In natural form,
\[
f(x;\eta)
=\exp\{\eta T(x)+d_0(\eta)+S(x)\}I_A(x),
\]
so
\[
\ell''(\eta;x)=d_0''(\eta).
\]
Regularity gives
\[
I(\eta;X)=-\E_\eta[\ell''(\eta;X)]=-d_0''(\eta).
\]

In the original parameterization,
\[
\ell(\theta;x)=c(\theta)T(x)+d(\theta)+S(x).
\]
The mean-zero score identity gives
\[
c'(\theta)\E_\theta[T(X)]+d'(\theta)=0,
\]
so, when $c'(\theta)\ne0$,
\[
\E_\theta[T(X)]=-\frac{d'(\theta)}{c'(\theta)}.
\]
Since
\[
\ell''(\theta;x)=c''(\theta)T(x)+d''(\theta),
\]
we obtain
\[
I(\theta;X)
=-\E_\theta[\ell''(\theta;X)]
=-d''(\theta)+\frac{d'(\theta)c''(\theta)}{c'(\theta)}.
\]
\end{solution}

\question[0] \textbf{Problem 1.14.}

Let $X_1,\ldots,X_n$ be iid $\mathcal{N}(\theta,\theta)$, $\theta>0$.
\begin{parts}
\part Find the Fisher information.
\part Find the efficiencies of $\bar X$ and $s^2$ and determine which is better.
\part Determine whether an unbiased estimator uniformly improves both.
\part Determine whether an unbiased estimator achieves efficiency one.
\end{parts}

\begin{solution}
Let
\[
s^2
=
\frac{1}{n-1}
\sum_{i=1}^n(X_i-\bar X)^2.
\]
Since $X_i\sim\mathcal{N}(\theta,\theta)$, both
\[
\E_\theta[\bar X]=\theta
\qquad\text{and}\qquad
\E_\theta[s^2]=\theta.
\]
Thus $\bar X$ and $s^2$ are both unbiased estimators of $\theta$.

\begin{parts}

\part
For one observation,
\[
X\sim\mathcal{N}(\mu(\theta),v(\theta)),
\]
where
\[
\mu(\theta)=\theta,
\qquad
v(\theta)=\theta.
\]
For a normal distribution with both mean and variance depending on
$\theta$,
\[
I_1(\theta)
=
\frac{\{\mu'(\theta)\}^2}{v(\theta)}
+
\frac{1}{2}
\frac{\{v'(\theta)\}^2}{v(\theta)^2}.
\]
Here,
\[
\mu'(\theta)=1,
\qquad
v'(\theta)=1,
\]
so
\[
I_1(\theta)
=
\frac{1}{\theta}
+
\frac{1}{2\theta^2}
=
\frac{2\theta+1}{2\theta^2}.
\]
Because the observations are independent,
\[
I_n(\theta)
=
nI_1(\theta)
=
\frac{n(2\theta+1)}{2\theta^2}.
\]
Therefore, the Cram\'er--Rao lower bound for any unbiased estimator of
$\theta$ is
\[
\frac{1}{I_n(\theta)}
=
\frac{2\theta^2}{n(2\theta+1)}.
\]

\part
We have
\[
\Var_\theta(\bar X)=\frac{\theta}{n},
\]
and, because $(n-1)s^2/\theta\sim\chi^2_{n-1}$,
\[
\Var_\theta(s^2)=\frac{2\theta^2}{n-1}.
\]

Define the efficiency of an unbiased estimator $\delta$ of $\theta$ by
\[
\operatorname{Eff}(\delta)
=
\frac{1/I_n(\theta)}{\Var_\theta(\delta)}.
\]
Then
\[
\operatorname{Eff}(\bar X)
=
\frac{\dfrac{2\theta^2}{n(2\theta+1)}}
     {\dfrac{\theta}{n}}
=
\frac{2\theta}{2\theta+1},
\]
and
\[
\operatorname{Eff}(s^2)
=
\frac{\dfrac{2\theta^2}{n(2\theta+1)}}
     {\dfrac{2\theta^2}{n-1}}
=
\frac{n-1}{n(2\theta+1)}.
\]

Thus,
\[
\operatorname{Eff}(\bar X)
>
\operatorname{Eff}(s^2)
\]
if and only if
\[
\frac{2\theta}{2\theta+1}
>
\frac{n-1}{n(2\theta+1)}.
\]
Since $2\theta+1>0$, this is equivalent to
\[
2n\theta>n-1,
\]
or
\[
\theta>\frac{n-1}{2n}.
\]

Therefore:
\[
\begin{cases}
\bar X \text{ is more efficient}, &
\theta>\dfrac{n-1}{2n},\\[6pt]
s^2 \text{ is more efficient}, &
0<\theta<\dfrac{n-1}{2n},\\[6pt]
\bar X \text{ and }s^2\text{ have equal efficiency}, &
\theta=\dfrac{n-1}{2n}.
\end{cases}
\]

\part
The joint density is
\begin{align*}
f_\theta(x_1,\ldots,x_n)
&=
\prod_{i=1}^n
\frac{1}{\sqrt{2\pi\theta}}
\exp\left\{
-\frac{(x_i-\theta)^2}{2\theta}
\right\}
\\
&=
(2\pi\theta)^{-n/2}
\exp\left\{
-\frac{1}{2\theta}\sum_{i=1}^n x_i^2
+\sum_{i=1}^n x_i
-\frac{n\theta}{2}
\right\}.
\end{align*}
Thus the joint density factors through
\[
T=\sum_{i=1}^nX_i^2.
\]
Hence $T$ is sufficient.

Moreover, this is a one-dimensional exponential family with natural statistic $T$
and natural parameter
\[
\eta=-\frac{1}{2\theta},
\qquad
\eta\in(-\infty,0).
\]
The representation is minimal (the single statistic $T$ is non-constant on the
support, so condition (b) of Problem 1.4 holds trivially with $k=1$), and the
natural parameter space $(-\infty,0)$ is open in $\R$. By the completeness
theorem for full-rank (minimal) exponential families with open natural parameter
space, $T$ is therefore complete sufficient.

Now define
\[
\delta(T)
=
\E_\theta[\bar X\mid T].
\]
Since $T$ is sufficient, the conditional distribution
$P_\theta(X\in\cdot\mid T=t)$ does not depend on $\theta$; write this common
conditional law as $P(\cdot\mid t)$. Defining $\delta(t)=\E[\bar X\mid T=t]$
using this single $\theta$-free kernel produces one function of $T$ alone
that is simultaneously a version of $\E_\theta[\bar X\mid T]$ for every
$\theta$. Since $\bar X$ is unbiased,
\[
\E_\theta[\delta(T)]
=
\E_\theta[\bar X]
=
\theta.
\]
Thus $\delta(T)$ is an unbiased function of the complete sufficient
statistic $T$, and hence it is the UMVUE of $\theta$.

Likewise,
\[
\E_\theta[s^2\mid T]
\]
is an unbiased function of $T$. Since both
\[
\E_\theta[\bar X\mid T]
\qquad\text{and}\qquad
\E_\theta[s^2\mid T]
\]
are unbiased functions of the complete sufficient statistic $T$,
completeness implies that they are equal almost surely. Therefore the
same estimator $\delta(T)$ is the Rao--Blackwell improvement of both
$\bar X$ and $s^2$.

Consequently,
\[
\Var_\theta(\delta(T))
\le
\Var_\theta(\bar X)
\]
and
\[
\Var_\theta(\delta(T))
\le
\Var_\theta(s^2)
\]
for every $\theta>0$. Hence an unbiased estimator uniformly improves
both $\bar X$ and $s^2$.

\part
The log likelihood, up to terms not depending on $\theta$, is
\[
\ell(\theta)
=
-\frac{n}{2}\log\theta
-\frac{T}{2\theta}
-\frac{n\theta}{2}.
\]
The score is
\[
U_\theta(X)
=
\frac{\partial\ell}{\partial\theta}
=
-\frac{n}{2\theta}
+
\frac{T}{2\theta^2}
-
\frac{n}{2}
=
\frac{T-n\theta-n\theta^2}{2\theta^2}.
\]

Suppose that an unbiased estimator $\delta(X)$ attained the
Cram\'er--Rao lower bound for every $\theta>0$. Equality in the
Cram\'er--Rao inequality would imply
\[
\delta(X)-\theta
=
\frac{1}{I_n(\theta)}U_\theta(X).
\]
Since
\[
\frac{1}{I_n(\theta)}
=
\frac{2\theta^2}{n(2\theta+1)},
\]
we would obtain
\begin{align*}
\delta(X)-\theta
&=
\frac{2\theta^2}{n(2\theta+1)}
\cdot
\frac{T-n\theta-n\theta^2}{2\theta^2}
\\
&=
\frac{T-n\theta-n\theta^2}{n(2\theta+1)}.
\end{align*}
Therefore,
\[
\delta(X)
=
\theta+
\frac{T-n\theta-n\theta^2}{n(2\theta+1)}
=
\frac{T+n\theta^2}{n(2\theta+1)}.
\]

The right-hand side depends nontrivially on the unknown parameter
$\theta$. Hence it cannot equal one fixed statistic $\delta(X)$ for
every $\theta>0$. Therefore no unbiased estimator attains efficiency
one uniformly over the parameter space.

\end{parts}
\end{solution}

\question[0] \textbf{Problem 1.15.}

Let $X_i$ be iid Poisson$(\lambda)$ and let $\widehat\gamma$ be the UMVUE of
$\gamma=e^\lambda$.
\begin{parts}
\part Find its mean and variance.
\part Show that $\Var(\widehat\gamma)>1/I(\gamma;X)$.
\part Show that the ratio tends to one as $n\to\infty$.
\end{parts}

\begin{solution}
Let
\[
T=\sum_{i=1}^nX_i\sim\operatorname{Poisson}(n\lambda).
\]
Using the Poisson pgf,
\[
\widehat\gamma=\left(1+\frac1n\right)^T
\]
is unbiased for $e^\lambda$. Since it is a function of the complete sufficient
statistic $T$, it is the UMVUE.

\begin{parts}
\part
The pgf gives
\[
\E[\widehat\gamma]=e^\lambda
\]
and
\[
\E[\widehat\gamma^2]
=\exp\left\{n\lambda\left(\left(1+\frac1n\right)^2-1\right)\right\}
=e^{2\lambda+\lambda/n}.
\]
Therefore
\[
\Var(\widehat\gamma)
=e^{2\lambda}\left(e^{\lambda/n}-1\right).
\]

\part
Since $I(\lambda;X)=n/\lambda$ and $\gamma=e^\lambda$, the reparametrization formula
gives
\[
I(\gamma;X)=\frac{n}{\lambda e^{2\lambda}},
\qquad
\frac1{I(\gamma;X)}=\frac{\lambda e^{2\lambda}}n.
\]
Because $e^x-1>x$ for $x>0$,
\[
e^{2\lambda}\left(e^{\lambda/n}-1\right)
>\frac{\lambda e^{2\lambda}}n.
\]
Hence $\Var(\widehat\gamma)>1/I(\gamma;X)$.

\part
The ratio is
\[
\frac{\Var(\widehat\gamma)}{1/I(\gamma;X)}
=\frac{e^{\lambda/n}-1}{\lambda/n}
\longrightarrow1
\qquad(n\to\infty).
\]
\end{parts}
\end{solution}

\question[0] \textbf{Problem 1.16.} (Optional)

If $T(X)\in\R^k$ is sufficient and there exist
$\theta_0,\ldots,\theta_k$ such that
\[
\left(
\frac{f(x;\theta_1)}{f(x;\theta_0)},\ldots,
\frac{f(x;\theta_k)}{f(x;\theta_0)}
\right)
\]
is a one-to-one function of $T(x)$, prove that $T$ is minimal sufficient.

\begin{solution}
Work on the common support where the displayed ratios are defined. Since $T$ is
sufficient, the factorization theorem implies that
\[
T(x)=T(y)
\quad\Longrightarrow\quad
\frac{f(x;\theta)}{f(y;\theta)}
\text{ is independent of }\theta.
\]
Conversely, suppose $f(x;\theta)/f(y;\theta)$ is independent of $\theta$. Then for
$i=1,\ldots,k$,
\[
\frac{f(x;\theta_i)/f(x;\theta_0)}
{f(y;\theta_i)/f(y;\theta_0)}=1.
\]
Hence the two ratio vectors at $x$ and $y$ are equal. By the assumed one-to-one
relationship with $T$, this implies $T(x)=T(y)$. Therefore
\[
T(x)=T(y)
\quad\Longleftrightarrow\quad
\frac{f(x;\theta)}{f(y;\theta)}
\text{ is independent of }\theta.
\]
The likelihood-ratio characterization of minimal sufficiency now shows that $T$ is
minimal sufficient.
\end{solution}

\question[0] \textbf{Problem 1.17.}

If $X\sim\operatorname{Gamma}(\alpha,\beta)$ in the shape--scale parameterization,
find $\E_{\alpha,\beta}[\log X]$, $\Var_{\alpha,\beta}(\log X)$, and
$\Cov_{\alpha,\beta}(\log X,X)$ using the exponential-family representation.

\begin{solution}
The density is
\[
f(x;\alpha,\beta)
=
\exp\left\{
(\alpha-1)\log x-\frac{x}{\beta}
-\log\Gamma(\alpha)-\alpha\log\beta
\right\}I_{(0,\infty)}(x).
\]
Take
\[
T_1(X)=\log X,
\qquad
T_2(X)=X,
\qquad
\eta_1=\alpha-1,
\qquad
\eta_2=-\frac1\beta.
\]
The log-partition function is
\[
A(\eta)
=\log\Gamma(\eta_1+1)-(\eta_1+1)\log(-\eta_2).
\]
Using $\nabla A(\eta)=\E_\eta[T(X)]$ (and $\E_\eta[T(X)]=\E_{\alpha,\beta}[T(X)]$
under the change of variables $\eta\leftrightarrow(\alpha,\beta)$),
\[
\E_{\alpha,\beta}[\log X]
=\frac{\partial A}{\partial\eta_1}
=\psi(\eta_1+1)-\log(-\eta_2)
=\psi(\alpha)+\log\beta.
\]
Using $\nabla^2A(\eta)=\Cov_\eta(T(X))=\Cov_{\alpha,\beta}(T(X))$,
\[
\Var_{\alpha,\beta}(\log X)
=\frac{\partial^2A}{\partial\eta_1^2}
=\psi_1(\alpha)
\]
and
\[
\Cov_{\alpha,\beta}(\log X,X)
=\frac{\partial^2A}{\partial\eta_1\partial\eta_2}
=-\frac1{\eta_2}
=\beta.
\]
\end{solution}

\question[0] \textbf{Problem 1.18.}

Let $X=(X_1,\ldots,X_n)$ be iid
$\operatorname{Gamma}(\alpha,\beta)$. Prove or disprove: if
\[
\E_{\alpha,\beta}[g(X)]=0
\]
for all $\alpha,\beta>0$, then
\[
\E_{\alpha,\beta}\left[
g(X)
\ \middle|\
\sum_{i=1}^n X_i,\,
\prod_{i=1}^n X_i
\right]=0
\]
for all $\alpha,\beta>0$.

\begin{solution}
The assertion is true, assuming $g(X)$ is integrable, i.e.\
$\E_{\alpha,\beta}|g(X)|<\infty$, at every $(\alpha,\beta)$ with $\alpha,\beta>0$
(this is implicit in the hypothesis $\E_{\alpha,\beta}[g(X)]=0$ for all
$\alpha,\beta>0$, and is needed for the conditional expectations below to be
well defined).

The joint density is
\begin{align*}
f_{\alpha,\beta}(x)
&=
\prod_{i=1}^n
\frac{1}{\Gamma(\alpha)\beta^\alpha}
x_i^{\alpha-1}e^{-x_i/\beta}
I_{(0,\infty)}(x_i)
\\
&=
\exp\left\{
(\alpha-1)\sum_{i=1}^n\log x_i
-\frac{1}{\beta}\sum_{i=1}^n x_i
-n\log\Gamma(\alpha)
-n\alpha\log\beta
\right\}
I_{(0,\infty)^n}(x).
\end{align*}

Therefore,
\[
T_1(X)=\sum_{i=1}^nX_i,
\qquad
T_2(X)=\sum_{i=1}^n\log X_i
\]
is a sufficient statistic. Its natural parameters are
\[
\eta_1=\alpha-1,
\qquad
\eta_2=-\frac{1}{\beta},
\]
with natural parameter space
\[
(-1,\infty)\times(-\infty,0),
\]
which is open in $\R^2$.

Also, the representation is minimal: the statistics
$1$, $T_1$, and $T_2$ are linearly independent. Hence, by the
full-rank exponential-family completeness theorem,
\[
(T_1,T_2)
=
\left(
\sum_{i=1}^nX_i,\,
\sum_{i=1}^n\log X_i
\right)
\]
is complete and sufficient.

Since
\[
\sum_{i=1}^n\log X_i
=
\log\left(\prod_{i=1}^nX_i\right),
\]
the two statistics
\[
\left(
\sum_{i=1}^nX_i,\,
\sum_{i=1}^n\log X_i
\right)
\]
and
\[
\left(
\sum_{i=1}^nX_i,\,
\prod_{i=1}^nX_i
\right)
\]
generate the same sigma-field.

By sufficiency, there exists a parameter-free version
\[
h(T_1,T_2)
=
\E_{\alpha,\beta}[g(X)\mid T_1,T_2].
\]
By the tower property,
\[
\E_{\alpha,\beta}[h(T_1,T_2)]
=
\E_{\alpha,\beta}[g(X)]
=
0
\]
for every $\alpha,\beta>0$. Since $(T_1,T_2)$ is complete,
\[
h(T_1,T_2)=0
\]
almost surely for every $\alpha,\beta>0$.

Therefore,
\[
\E_{\alpha,\beta}\left[
g(X)
\ \middle|\
\sum_{i=1}^nX_i,\,
\prod_{i=1}^nX_i
\right]
=
0
\]
almost surely for all $\alpha,\beta>0$.
\end{solution}

\question[0] \textbf{Problem 1.19.}

Suppose $Y_1,\ldots,Y_n$ are independent with
\[
Y_i\sim\mathcal{N}(\beta x_i,\sigma^2),
\]
where the $x_i$ are fixed and known.
\begin{parts}
\part Find complete sufficient statistics.
\part Find the MLEs of $(\beta,\sigma)$.
\part Find the UMVUEs of $\beta$ and $\sigma^2$.
\end{parts}

\begin{solution}
Let
\[
q=\sum_{i=1}^n x_i^2>0,
\qquad
T_1=\sum_{i=1}^n x_iY_i,
\qquad
T_2=\sum_{i=1}^n Y_i^2.
\]

\begin{parts}
\part
Expanding the joint density gives
\[
f_{\beta,\sigma^2}(y)
=(2\pi)^{-n/2}
\exp\left\{
\frac{\beta}{\sigma^2}T_1
-\frac{1}{2\sigma^2}T_2
-\frac{\beta^2q}{2\sigma^2}
-\frac n2\log\sigma^2
\right\}.
\]
Thus $(T_1,T_2)$ is sufficient. The natural parameters are
\[
\eta_1=\frac{\beta}{\sigma^2},
\qquad
\eta_2=-\frac{1}{2\sigma^2},
\]
and they range over the open set $\R\times(-\infty,0)$.

The design is assumed nondegenerate, i.e.\ not every $x_i$ is $0$, so that
$q=\sum_{i=1}^nx_i^2>0$; this is exactly the condition needed for minimality. Indeed,
by the criterion of Problem 1.4, minimality requires that
$b_1T_1(y)+b_2T_2(y)=b_0$ for all $y\in\R^n$ forces $b_0=b_1=b_2=0$. Comparing the
coefficient of $y_i^2$ gives $b_2=0$; comparing the coefficient of $y_i$ then gives
$b_1x_i=0$ for every $i$, and since $q>0$ some $x_i\ne0$, so $b_1=0$; finally
$b_0=0$. Hence the representation is minimal (for every $n\ge1$; we additionally
take $n\ge2$ below so that the UMVUE of $\sigma^2$ in part (c) has positive degrees
of freedom). Hence the full-rank exponential-family completeness theorem shows that
$(T_1,T_2)$ is complete sufficient.

\part
Least squares gives
\[
\widehat\beta=\frac{T_1}{q}.
\]
The residual sum of squares is
\[
\operatorname{RSS}
=\sum_{i=1}^n(Y_i-\widehat\beta x_i)^2
=T_2-\frac{T_1^2}{q}.
\]
Therefore, when $\operatorname{RSS}>0$,
\[
\widehat{\sigma^2}_{\mathrm{MLE}}=\frac{\operatorname{RSS}}{n},
\qquad
\widehat\sigma_{\mathrm{MLE}}=\sqrt{\frac{\operatorname{RSS}}{n}}.
\]
If $\operatorname{RSS}=0$, the likelihood is unbounded as $\sigma^2\downarrow0$, so no
MLE exists.

\part
The estimator $\widehat\beta=T_1/q$ is unbiased and is a function of the complete
sufficient statistic. Therefore
\[
\widehat\beta_{\mathrm{UMVUE}}
=\frac{\sum_{i=1}^n x_iY_i}{\sum_{i=1}^n x_i^2}.
\]
For $n\ge2$,
\[
\frac{\operatorname{RSS}}{\sigma^2}\sim\chi^2_{n-1},
\]
so $\E[\operatorname{RSS}]=(n-1)\sigma^2$. Hence
\[
\widehat{\sigma^2}_{\mathrm{UMVUE}}
=\frac{\operatorname{RSS}}{n-1}
=\frac{1}{n-1}\left(T_2-\frac{T_1^2}{q}\right).
\]
Both estimators are functions of the complete sufficient statistic, so the
Lehmann--Scheff\'e theorem proves the UMVUE claims.
\end{parts}
\end{solution}

\end{questions}

\end{document}